\def\tituloconcepto{Algoritmo de Deutsch}
% Si el capitulo se compila solo (sin main.tex), esta cabecera abre el documento.
% Antes de % Si el capitulo se compila solo (sin main.tex), esta cabecera abre el documento.
% Antes de % Si el capitulo se compila solo (sin main.tex), esta cabecera abre el documento.
% Antes de \input{cabecera} cada capitulo define \tituloconcepto.
\ifdefined\enmain\else
  \documentclass[11pt,a4paper]{article}
  \input{preambulo}
  \begin{document}
  \title{\tituloconcepto}
  \author{\autor}
  \date{\today}
  \maketitle
\fi
 cada capitulo define \tituloconcepto.
\ifdefined\enmain\else
  \documentclass[11pt,a4paper]{article}
  % Preambulo compartido por todos los documentos. Compilar con XeLaTeX/LuaLaTeX
% (o Tectonic). Si se usa Overleaf, elegir el compilador XeLaTeX.
\usepackage[spanish,es-noshorthands]{babel}
\usepackage{amsmath,amssymb}
\usepackage{booktabs,array}
\usepackage[a4paper,margin=2.5cm]{geometry}
\usepackage{xcolor}
\usepackage{listings}
\usepackage[most]{tcolorbox}
\usepackage{quantikz}
\usepackage[hidelinks]{hyperref}

% ---- Autor (editar a gusto) -------------------------------------------------
\newcommand{\autor}{Genesis Silesky}

% ---- Notacion de Dirac y atajos ---------------------------------------------
% quantikz ya define \ket y \bra; DeclareRobustCommand los redefine sin error.
\DeclareRobustCommand{\ket}[1]{\left|#1\right\rangle}
\DeclareRobustCommand{\bra}[1]{\left\langle#1\right|}
\DeclareRobustCommand{\braket}[2]{\left\langle#1\middle|#2\right\rangle}
\newcommand{\C}{\mathbb{C}}
\newcommand{\R}{\mathbb{R}}
\DeclareMathOperator{\Tr}{Tr}

% ---- Cajas de apoyo --------------------------------------------------------
% "concepto": definiciones previas que hacen falta para entender el capitulo.
\newtcolorbox{concepto}[1]{breakable,colback=blue!4,colframe=blue!55!black,
  fonttitle=\bfseries,title={Concepto previo: #1}}
% "nota": aclaraciones y advertencias.
\newtcolorbox{nota}[1]{breakable,colback=orange!6,colframe=orange!70!black,
  fonttitle=\bfseries,title={#1}}

% ---- Estilos de listados ----------------------------------------------------
\lstdefinestyle{py}{language=Python,basicstyle=\ttfamily\footnotesize,
  breaklines=true,postbreak=\mbox{\textcolor{gray}{$\hookrightarrow$}\space},
  numbers=left,numberstyle=\tiny\color{gray},frame=single,
  commentstyle=\color{green!40!black}\itshape,
  keywordstyle=\color{blue!70!black}\bfseries,
  stringstyle=\color{red!60!black},showstringspaces=false,tabsize=4,
  columns=fullflexible,keepspaces=true}
\lstdefinestyle{salida}{basicstyle=\ttfamily\footnotesize,breaklines=true,
  frame=leftline,backgroundcolor=\color{gray!8},columns=fullflexible,
  keepspaces=true}

  \begin{document}
  \title{\tituloconcepto}
  \author{\autor}
  \date{\today}
  \maketitle
\fi
 cada capitulo define \tituloconcepto.
\ifdefined\enmain\else
  \documentclass[11pt,a4paper]{article}
  % Preambulo compartido por todos los documentos. Compilar con XeLaTeX/LuaLaTeX
% (o Tectonic). Si se usa Overleaf, elegir el compilador XeLaTeX.
\usepackage[spanish,es-noshorthands]{babel}
\usepackage{amsmath,amssymb}
\usepackage{booktabs,array}
\usepackage[a4paper,margin=2.5cm]{geometry}
\usepackage{xcolor}
\usepackage{listings}
\usepackage[most]{tcolorbox}
\usepackage{quantikz}
\usepackage[hidelinks]{hyperref}

% ---- Autor (editar a gusto) -------------------------------------------------
\newcommand{\autor}{Genesis Silesky}

% ---- Notacion de Dirac y atajos ---------------------------------------------
% quantikz ya define \ket y \bra; DeclareRobustCommand los redefine sin error.
\DeclareRobustCommand{\ket}[1]{\left|#1\right\rangle}
\DeclareRobustCommand{\bra}[1]{\left\langle#1\right|}
\DeclareRobustCommand{\braket}[2]{\left\langle#1\middle|#2\right\rangle}
\newcommand{\C}{\mathbb{C}}
\newcommand{\R}{\mathbb{R}}
\DeclareMathOperator{\Tr}{Tr}

% ---- Cajas de apoyo --------------------------------------------------------
% "concepto": definiciones previas que hacen falta para entender el capitulo.
\newtcolorbox{concepto}[1]{breakable,colback=blue!4,colframe=blue!55!black,
  fonttitle=\bfseries,title={Concepto previo: #1}}
% "nota": aclaraciones y advertencias.
\newtcolorbox{nota}[1]{breakable,colback=orange!6,colframe=orange!70!black,
  fonttitle=\bfseries,title={#1}}

% ---- Estilos de listados ----------------------------------------------------
\lstdefinestyle{py}{language=Python,basicstyle=\ttfamily\footnotesize,
  breaklines=true,postbreak=\mbox{\textcolor{gray}{$\hookrightarrow$}\space},
  numbers=left,numberstyle=\tiny\color{gray},frame=single,
  commentstyle=\color{green!40!black}\itshape,
  keywordstyle=\color{blue!70!black}\bfseries,
  stringstyle=\color{red!60!black},showstringspaces=false,tabsize=4,
  columns=fullflexible,keepspaces=true}
\lstdefinestyle{salida}{basicstyle=\ttfamily\footnotesize,breaklines=true,
  frame=leftline,backgroundcolor=\color{gray!8},columns=fullflexible,
  keepspaces=true}

  \begin{document}
  \title{\tituloconcepto}
  \author{\autor}
  \date{\today}
  \maketitle
\fi


\section{Algoritmo de Deutsch}

\subsection{Objetivo}
Resolver con una sola consulta un problema que clásicamente necesita dos:
decidir si una función de un bit es constante o balanceada.

\subsection{Conceptos previos}
\begin{concepto}{superposición, puerta H e interferencia}
Ver capítulos~01, 02 y~05a. La puerta $H$ convierte $\ket0$ en
$(\ket0+\ket1)/\sqrt2$ y $\ket1$ en $(\ket0-\ket1)/\sqrt2$, y dos $H$ seguidas
se cancelan.
\end{concepto}

\subsection{Fundamento teórico}
El algoritmo de Deutsch, propuesto por David Deutsch en 1985, se considera el
primer algoritmo que demuestra la ventaja de un computador cuántico sobre uno
clásico.

\paragraph{El problema (las cajas de pintura).} Hay dos bolas, una liviana
($x=0$) y una pesada ($x=1$), que se pintan de blanco ($f=0$) o negro ($f=1$).
Hay cajas de un solo color (pintan las dos bolas igual) y cajas de dos colores
(las pintan distinto), pero se mezclaron. ¿Cuántas veces hay que pintar para
distinguirlas?

En términos matemáticos, la caja es una función $f:\{0,1\}\to\{0,1\}$:
\begin{itemize}
  \item \textbf{constante} si $f(0)=f(1)$ (caja de un solo color);
  \item \textbf{balanceada} si $f(0)\neq f(1)$ (caja de dos colores).
\end{itemize}
Clásicamente hay que evaluar $f(0)$ y $f(1)$ y compararlos: \textbf{dos consultas}.

\paragraph{El oráculo.} La función se da como una puerta unitaria $U_f$ que
actúa como
\begin{equation}
  U_f\ket x\ket y=\ket x\ket{y\oplus f(x)},
\end{equation}
donde $\oplus$ es la suma módulo~2 (XOR). El primer qubit es el
\emph{registro de datos} y el segundo el \emph{registro objetivo}.
Gracias a la superposición, una sola aplicación de $U_f$ evalúa $f(0)$ y $f(1)$
a la vez (\emph{paralelismo cuántico}).

\paragraph{El circuito.}
\begin{center}
\begin{quantikz}
\lstick{$\ket0$} & \gate{H} & \gate[2]{U_f} & \gate{H} & \meter{}\\
\lstick{$\ket1$} & \gate{H} &               & \qw      & \qw
\end{quantikz}
\end{center}
\begin{enumerate}
  \item Estado inicial $\ket{\psi_0}=\ket0\ket1$.
  \item Con $H$ en ambos qubits:
    $\ket{\psi_1}=\frac{\ket0+\ket1}{\sqrt2}\,\frac{\ket0-\ket1}{\sqrt2}$.
  \item Una consulta a $U_f$. Como el objetivo está en $(\ket0-\ket1)/\sqrt2$,
    el oráculo solo cambia el signo de cada $\ket x$:
    \begin{equation}
      U_f\ket x\,\frac{\ket0-\ket1}{\sqrt2}=(-1)^{f(x)}\ket x\,\frac{\ket0-\ket1}{\sqrt2}.
    \end{equation}
    Por tanto
    $\ket{\psi_2}=\frac{(-1)^{f(0)}\ket0+(-1)^{f(1)}\ket1}{\sqrt2}\,
    \frac{\ket0-\ket1}{\sqrt2}$.
    \begin{itemize}
      \item Si $f$ es constante, los datos quedan en $\pm\frac{\ket0+\ket1}{\sqrt2}$.
      \item Si $f$ es balanceada, quedan en $\pm\frac{\ket0-\ket1}{\sqrt2}$.
    \end{itemize}
  \item $H$ sobre los datos: $\ket{\psi_3}=\pm\ket0$ si $f$ es constante y
    $\pm\ket1$ si es balanceada.
  \item Se mide el registro de datos: $0$ significa constante y $1$ significa balanceada.
\end{enumerate}
\textbf{Una consulta contra dos.}

\subsection{Código (completamente comentado)}
\lstinputlisting[style=py]{../../fuentes/15_deutsch.py}

\subsection{Explicación celda por celda}
\begin{description}
  \item[Celda 2] Lista las cuatro funciones posibles y dice si cada una es
    constante o balanceada.
  \item[Celda 3] Construye el oráculo $U_f$: una $X$ al objetivo si $f$ vale 1
    siempre, una CNOT si $f(x)=x$, y CNOT más $X$ si $f(x)=1-x$.
  \item[Celda 4] Arma el circuito: objetivo en $\ket1$, $H$ en ambos, oráculo,
    $H$ en los datos y medición.
  \item[Celda 5] Lo ejecuta con las cuatro funciones y 1000 mediciones.
  \item[Celda 6] Muestra los signos $(-1)^{f(x)}$ del registro de datos
    después del oráculo.
\end{description}

\subsection{Salida obtenida}
\lstinputlisting[style=salida]{../../fuentes/salidas/15_deutsch.txt}
(El dibujo del circuito no se reproduce aquí.)

\subsection{Interpretación}
El resultado no es aleatorio: las dos funciones constantes siempre dan $0$ y las
dos balanceadas siempre dan $1$, con una sola consulta. Los signos de la última
parte coinciden con $(-1)^{f(x)}$: $(+1,+1)$ y $(-1,-1)$ para las constantes,
$(+1,-1)$ y $(-1,+1)$ para las balanceadas. Lo que distingue los casos es la
\emph{fase relativa} entre $\ket0$ y $\ket1$, que la $H$ final convierte en
una diferencia medible (interferencia, capítulo~05a).

\subsection{Ejercicios}
\begin{enumerate}
  \item Elimina la $X$ inicial del objetivo y mira qué cambia en la
    Celda~6.
  \item ¿Cuántas consultas necesitaría un algoritmo clásico para $n=1$? ¿Y
    el cuántico?
  \item Cambia el oráculo de $f(x)=x$ por otro circuito equivalente y comprueba que
    el resultado no cambia.
\end{enumerate}

\subsection{Errores comunes}
\begin{itemize}
  \item Creer que el algoritmo conoce $f(0)$ y $f(1)$ por separado: solo
    averigua si son iguales o distintos.
  \item Olvidar poner el objetivo en $\ket1$ antes de la $H$: sin esa fase
    el oráculo no cambia el signo.
\end{itemize}

\ifdefined\enmain\else
  \end{document}
\fi

